\documentclass{beamer}

% Packages
\usepackage{booktabs}
\usepackage{tabularx}
\usepackage{pifont}
\usepackage{microtype}
\usepackage{xcolor}

% Clean minimalist white theme
\usetheme{default}
\setbeamertemplate{navigation symbols}{}
\setbeamercolor{background canvas}{bg=white}
\setbeamercolor{frametitle}{fg=black, bg=white}
\setbeamercolor{normal text}{fg=black}
\setbeamertemplate{frametitle}{%
    \vspace{0.4cm}\textbf{\insertframetitle}\par\vskip-0.15cm\hrulefill}

% Check / cross marks
\newcommand{\cmark}{\textcolor{green!60!black}{\ding{51}}}
\newcommand{\xmark}{\textcolor{red}{\ding{55}}}

\begin{document}

\begin{frame}{Problem 9 [MEDQA]: Question}
    \textbf{Question:}\\[4pt]
    \small A 23-year-old woman comes to the physician because she is embarrassed about the appearance
of her nails. She has no history of serious illness and takes no medications. She appears
well. A photograph of the nails is shown. Which of the following additional findings is
most likely in this patient?\\[6pt]
    \vspace{2pt}\textbf{Options:}\\[2pt]
    \begin{description}
      \item[A.] Silvery plaques on extensor surfaces
      \item[B.] Flesh-colored papules in the lumbosacral region
      \item[C.] Erosions of the dental enamel
      \item[D.] Pallor of the conjunctival mucosa
      \item[E.] Holosystolic murmur at the left lower sternal border
    \end{description}
    \vspace{6pt}
    \textbf{True:} A \quad \textbf{Predicted:} A \quad \cmark
\end{frame}

\begin{frame}[allowframebreaks]{Problem 9: Reasoning (raw\_model (0))}
    \small The patient is a 23-year-old woman who is embarrassed about the appearance of her nails. The key detail here is that the physician has been shown a photograph of the nails. This suggests that the condition is related to nail appearance rather than systemic symptoms or other physical findings.

\medskip
Option A: Silvery plaques on extensor surfaces - This could refer to the characteristic appearance of psoriasis, which often presents as silvery scales on the skin, including the extensor surfaces of the limbs. Psoriasis can also affect the nails, causing changes such as pitting, discoloration, or lifting of the nail plate from the nail bed. This is a plausible association with nail appearance.

\medskip
Option B: Flesh-colored papules in the lumbosacral region - This finding is more likely to be associated with conditions like molluscum contagiosum or other dermatological conditions affecting the skin, not specifically nail changes.

\medskip
Option C: Erosions of the dental enamel - This would be more indicative of a systemic condition affecting the teeth and would not typically be associated with nail changes.

\medskip
Option D: Pallor of the conjunctival mucosa - This finding is more indicative of anemia or other systemic conditions affecting blood perfusion, not specifically nail changes.

\medskip
Option E: Holosystolic murmur at the left lower sternal border - This finding is more indicative of a cardiac condition and would not typically be associated with nail changes.

\medskip
Given the context of nail appearance and the options provided, the most likely additional finding related to the patient's nail condition would be silvery plaques on extensor surfaces, which is characteristic of psoriasis.

\medskip
Therefore, the most likely additional finding is:
\end{frame}

\begin{frame}[allowframebreaks]{Problem 9: Observation Coverage (raw\_model (0))}
    \begin{center}
    \textbf{Coverage score: } {\Large\textbf{2/10 = 20.0\%}}
    \end{center}
    \vspace{4pt}
    \textbf{Overall analysis:}\\[4pt]
    \small The chosen hypothesis (A: silvery plaques on extensor surfaces / psoriasis) accounts
for the explicit link between nail-related findings and psoriasis features, including
nails changes. Several raw details about the patient and the other option
descriptions are not explicitly connected to this hypothesis within the trace, so
they remain unaddressed. The trace does provide a direct, explicit connection for the
key detail about psoriasis features and the final conclusion, but it does not
logically tie the other enumerated details to psoriasis; thus, some supporting
observation details are not explained by the hypothesis.
    \vspace{8pt}
    \textbf{Observation details:}\\[4pt]
    \begin{tabularx}{\textwidth}{>{\centering\arraybackslash}p{0.9cm} X p{4.5cm}}
    \toprule
    \textbf{OK?} & \textbf{Detail} & \textbf{Evidence} \\
    \midrule
      \xmark & \small The patient is a 23-year-old woman & \small\itshape  \\
      \xmark & \small Embarrassed about the appearance of her nails & \small\itshape  \\
      \xmark & \small The physician has been shown a photograph of the nails & \small\itshape  \\
      \xmark & \small This suggests that the condition is related to nail
appearance rather than systemic symptoms or other physical
findings & \small\itshape  \\
      \cmark & \small Silvery plaques on extensor surfaces & \small\itshape "This is a plausible association with
nail appearance." \\
      \xmark & \small Flesh-colored papules in the lumbosacral region & \small\itshape  \\
      \xmark & \small Erosions of the dental enamel & \small\itshape  \\
      \xmark & \small Pallor of the conjunctival mucosa & \small\itshape  \\
      \xmark & \small Holosystolic murmur at the left lower sternal border & \small\itshape  \\
      \cmark & \small Conclusion that the most likely additional finding related
to the nail condition would be silvery plaques on extensor
surfaces, which is characteristic of psoriasis & \small\itshape "which is characteristic of psoriasis." \\
    \bottomrule
    \end{tabularx}
\end{frame}

\begin{frame}[allowframebreaks]{Problem 9: Reasoning (checkpoint-2560)}
    \small The patient's complaint is about the appearance of her nails, which suggests a dermatological issue. The options provided include various skin conditions and systemic manifestations. 

\medskip
Option A: Silvery plaques on extensor surfaces could refer to psoriasis, which often presents as silvery scales on the elbows, knees, and other extensor surfaces. This is a common dermatological condition that can affect nail appearance.

\medskip
Option B: Flesh-colored papules in the lumbosacral region typically refer to seborrheic keratoses or other benign skin growths, which are not directly related to nail changes.

\medskip
Option C: Erosions of the dental enamel would be more indicative of a systemic condition affecting oral health rather than nail appearance.

\medskip
Option D: Pallor of the conjunctival mucosa might indicate anemia but does not directly relate to nail appearance.

\medskip
Option E: A holosystolic murmur at the left lower sternal border would suggest a cardiac issue, which is unrelated to nail appearance.

\medskip
Given the context of the patient's complaint about nail appearance, the most likely additional finding would be silvery plaques on extensor surfaces, which is characteristic of psoriasis.
\end{frame}

\begin{frame}[allowframebreaks]{Problem 9: Observation Coverage (checkpoint-2560)}
    \begin{center}
    \textbf{Coverage score: } {\Large\textbf{8/9 = 88.9\%}}
    \end{center}
    \vspace{4pt}
    \textbf{Overall analysis:}\\[4pt]
    \small The chosen hypothesis (psoriasis with silvery plaques on extensor surfaces)
adequately accounts for the key observation that nail appearance prompted a
dermatological interpretation and is directly linked to the most plausible associated
finding described in the trace. Several peripheral details (B–E) are acknowledged as
not directly related to nail changes, while the primary link between nail appearance
and psoriasis is explicitly stated. The trace connects the core observation to the
hypothesis, but some minor contextual options are not used to justify the hypothesis,
which is appropriate. Overall, the hypothesis is well-supported by the central
details but excludes unrelated data with no explanation required.
    \vspace{8pt}
    \textbf{Observation details:}\\[4pt]
    \begin{tabularx}{\textwidth}{>{\centering\arraybackslash}p{0.9cm} X p{4.5cm}}
    \toprule
    \textbf{OK?} & \textbf{Detail} & \textbf{Evidence} \\
    \midrule
      \cmark & \small The patient's complaint is about the appearance of her
nails. & \small\itshape "The patient's complaint is about the
appearance of her nails, which suggests
a dermatological issue." \\
      \cmark & \small This complaint suggests a dermatological issue. & \small\itshape "which suggests a dermatological issue." \\
      \xmark & \small Options include various skin conditions and systemic
manifestations. & \small\itshape  \\
      \cmark & \small Option A describes silvery plaques on extensor surfaces,
which could refer to psoriasis and can affect nail
appearance. & \small\itshape "Option A: Silvery plaques on extensor
surfaces could refer to psoriasis, which
often presents as silvery scales on the
elbows, knees, and other extensor
surfaces. This is a common
dermatological condition that can affect
nail appearance." \\
      \cmark & \small Option B describes flesh-colored papules in the lumbosacral
region typically referring to seborrheic keratoses or benign
skin growths, not directly related to nail changes. & \small\itshape "Option B: Flesh-colored papules in the
lumbosacral region typically refer to
seborrheic keratoses or other benign
skin growths, which are not directly
related to nail changes." \\
      \cmark & \small Option C describes erosions of the dental enamel, indicating
a systemic condition affecting oral health rather than nail
appearance. & \small\itshape "Option C: Erosions of the dental enamel
would be more indicative of a systemic
condition affecting oral health rather
than nail appearance." \\
      \cmark & \small Option D describes pallor of the conjunctival mucosa, which
might indicate anemia but does not directly relate to nail
appearance. & \small\itshape "Option D: Pallor of the conjunctival
mucosa might indicate anemia but does
not directly relate to nail appearance." \\
      \cmark & \small Option E describes a holosystolic murmur at the left lower
sternal border, suggesting a cardiac issue unrelated to nail
appearance. & \small\itshape "Option E: A holosystolic murmur at the
left lower sternal border would suggest
a cardiac issue, which is unrelated to
nail appearance." \\
      \cmark & \small Given the context, the most likely additional finding would
be silvery plaques on extensor surfaces, characteristic of
psoriasis. & \small\itshape "Given the context of the patient's
complaint about nail appearance, the
most likely additional finding would be
silvery plaques on extensor surfaces,
which is characteristic of psoriasis." \\
    \bottomrule
    \end{tabularx}
\end{frame}

\end{document}
